\chapter{Installation}

\section{Hardware}

\begin{itemize}
  \item \textbf{PolarFire SoC Discovery Kit} board (MPFS095T).
  \item microSD card with the board's Linux image.
  \item Ethernet cable to the client PC's network.
  \item USB cable, only for programming the bitstream (JTAG) and for the
    serial console.
\end{itemize}

A delivered device already has the bitstream programmed and the software
installed: it only needs connecting to the network. The following sections
describe an installation from scratch.

\section{Programming the bitstream}

The bitstream is programmed over JTAG from a PC with Microchip Libero SoC or
FlashPro Express, with the board's embedded programmer connected over USB. The
FPGA keeps it without power.

\begin{lstlisting}[language=shell]
# with FlashPro Express and the release's programming file
FPExpress  # open mpeg2fpga-<version>.job and run PROGRAM
\end{lstlisting}

Once Linux is up, the driver reports which bitstream is loaded:

\begin{lstlisting}[language=shell]
ssh root@board cat /sys/bus/platform/drivers/mpeg2fpga/*/build
0.1.0+8004bc9
\end{lstlisting}

That is the release version and the commit it was synthesized from; it ends
in \cmd{-dirty} if it was synthesized with uncommitted changes.

\section{Board software}

First the driver and its device tree overlay, then the service:

\begin{lstlisting}[language=shell]
scp -r driver/mpeg2fpga/tools root@board:/tmp/m2f-driver
ssh root@board sh /tmp/m2f-driver/install-on-board.sh

scp -r daemon api/PROTOCOL-v1.md root@board:/tmp/m2f-daemon/
ssh root@board sh /tmp/m2f-daemon/daemon/install-on-board.sh
\end{lstlisting}

The service installer:
\begin{itemize}
  \item creates the \cmd{mpeg2fpgad} system user (unprivileged);
  \item installs the service in \cmd{/usr/local/lib/mpeg2fpgad} and the native
    converter \cmd{libm2fconv.so};
  \item generates the access token in \cmd{/etc/mpeg2fpgad/token};
  \item enables and starts the systemd unit \cmd{mpeg2fpgad.service}.
\end{itemize}

Options: \cmd{-{}-tls} for HTTPS and \cmd{-{}-debug} for the diagnostic
endpoints (see chapter~\ref{cap:security}).

\section{Checking}

\begin{lstlisting}[language=shell]
curl http://192.168.18.5:8080/v1/health
{"ok": true}
\end{lstlisting}

\section{Client library}

On the PC, with Python 3.9 or later:

\begin{lstlisting}[language=shell]
pip install mpeg2fpga-<version>-py3-none-any.whl
ssh root@board cat /etc/mpeg2fpgad/token > token
\end{lstlisting}
